\chapter{Auslander--Buchsbaum}
\label{chap:Albanese_AuslanderBuchsbaum}
% archon:covers AlgebraicJacobian/Albanese/AuslanderBuchsbaum.lean

% STRATEGY NOTE (iter-174). This chapter is the A.4.b sub-row of Route A
% (positive-genus arm of nonempty_jacobianWitness). It packages the
% Auslander--Buchsbaum formula and the corollary ``regular local ring
% \(\Rightarrow\) Cohen--Macaulay'' that A.4.a (the codim-1 extension proof for
% Albanese morphisms out of a regular projective surface) consumes. Per
% STRATEGY.md L30 this row is gated downstream on A.4.a but is independently
% startable on the Mathlib side: the algebra here is decoupled from the
% Albanese geometry and lives entirely under \texttt{RingTheory.*}. The
% target Lean file \texttt{AlgebraicJacobian/Albanese/AuslanderBuchsbaum.lean}
% does not yet exist; this chapter is the prover-ready specification for
% creating it. Note that Mathlib already houses parts of the depth /
% projective-dimension API
% (\texttt{Mathlib.RingTheory.Depth.*},
% \texttt{Mathlib.RingTheory.Ideal.RegularSequence},
% \texttt{Mathlib.CategoryTheory.ProjectiveResolution}); a Plan-Agent audit
% at the pinned commit (\texttt{b80f227}) should confirm what is present
% before the prover writes the file, since a fully Mathlib-backed
% Auslander--Buchsbaum would collapse this chapter to a re-export.

\section{Setup and motivation}
\label{sec:ab_setup}

A regular projective surface \(S\) over a field --- such as \(\mathbb{P}^1 \times \mathbb{P}^1\)
or the product \(C \times C\) of a smooth projective curve \(C\) with itself --- has a regular
local ring at every (closed) point, of Krull dimension \(2\). To extend a rational map
\(f : S \dashrightarrow A\) from \(S\) to an abelian variety across a codimension-\(2\) closed
point, one argues that the local ring \(\mathcal{O}_{S,x}\) at such a point has \emph{depth at
least} \(2\), hence a rational function with no codimension-\(1\) poles on a neighbourhood of \(x\)
extends across \(x\). The depth \(\geq 2\) inequality is exactly what the
Auslander--Buchsbaum formula gives: for a regular local ring \(R\) one has \(\text{pd}_R(R) = 0\),
so
\[
  \text{depth}(R) \;=\; \text{pd}_R(R) + \text{depth}(R) \;=\; \text{depth}(R)
\]
trivially, but combined with the regular-sequence definition of depth and the fact that a
regular local ring of dimension \(d\) admits an \(R\)-regular sequence of length \(d\) generating
the maximal ideal, one obtains \(\text{depth}(R) = \dim(R)\), i.e.\ \(R\) is Cohen--Macaulay
(\cref{cor:regular_cohen_macaulay} below). This is the input A.4.a needs.

This chapter packages:
\begin{enumerate}
  \item the regular-sequence definition of depth and its \(\text{Ext}\) reformulation
    (\cref{def:depth} and \texttt{lem:depth_via_ext}),
  \item the projective-dimension invariant of a finitely-generated module
    (\texttt{def:projective_dimension}),
  \item the depth-on-a-short-exact-sequence lemma
    (\texttt{lem:depth_short_exact_sequence}),
  \item the Auslander--Buchsbaum formula proper
    (\cref{thm:auslander_buchsbaum}), and
  \item the corollary regular \(\Rightarrow\) Cohen--Macaulay
    (\cref{cor:regular_cohen_macaulay}).
\end{enumerate}
A short closing paragraph (\cref{sec:ab_application_to_a4a}) explains how A.4.a
consumes \cref{cor:regular_cohen_macaulay} to push a codim-1 extension across a
codim-2 point.

\section{Depth of a module over a local ring}
\label{sec:ab_depth}

\begin{definition}

\leanok
  [Depth of a module]
  \label{def:depth}
  \lean{RingTheory.Module.depth}
  % SOURCE: [Stacks Project], tag 00LF (definition-depth)
  % (read from references/stacks-algebra.tex, L17748-L17759).
  % SOURCE QUOTE:
  % "Let \(R\) be a ring, and \(I \subset R\) an ideal. Let \(M\) be a finite \(R\)-module.
  %  The {\it \(I\)-depth} of \(M\), denoted \(\text{depth}_I(M)\), is defined as follows:
  %  \begin{enumerate}
  %  \item if \(IM \not = M\), then \(\text{depth}_I(M)\) is the supremum in
  %  \(\{0, 1, 2, \ldots, \infty\}\) of the lengths of \(M\)-regular sequences in \(I\),
  %  \item if \(IM = M\) we set \(\text{depth}_I(M) = \infty\).
  %  \end{enumerate}
  %  If \((R, \mathfrak m)\) is local we call \(\text{depth}_{\mathfrak m}(M)\) simply
  %  the {\it depth} of \(M\)."
  \textit{Source: [Stacks Project], tag 00LF (definition-depth).}
  Let \(R\) be a commutative ring, let \(I \subseteq R\) be an ideal, and let \(M\) be a finitely
  generated \(R\)-module. The \emph{\(I\)-depth} of \(M\), written \(\text{depth}_I(M)\), is the
  supremum in \(\{0, 1, 2, \dots, \infty\}\) of the lengths of \(M\)-regular sequences contained
  in \(I\), provided \(IM \neq M\); if \(IM = M\) we set \(\text{depth}_I(M) = \infty\). When
  \((R, \mathfrak{m})\) is a Noetherian local ring we abbreviate
  \(\text{depth}(M) := \text{depth}_{\mathfrak{m}}(M)\) and call this simply the
  \emph{depth} of \(M\). The depth of the ring \(R\) itself is \(\text{depth}(R)\), the depth
  of \(R\) as a module over itself.
\end{definition}

Recall that a sequence \(f_1, \dots, f_r \in R\) is called \emph{\(M\)-regular} if
\(M/(f_1, \dots, f_r)M \neq 0\) and for each \(i\), the element \(f_i\) is a non-zero-divisor on
\(M/(f_1, \dots, f_{i-1})M\). The empty sequence is regular on \(M\) iff \(M \neq 0\) (it is
\emph{not} regular on the zero module, but the convention \(\text{depth}(0) = \infty\)
above is the usual one).

The two practical computations of depth are the regular-sequence definition (used to
establish a \emph{lower} bound) and the \(\text{Ext}\) characterisation (used to establish an
\emph{upper} bound and to prove the depth-on-a-short-exact-sequence inequality).


% SOURCE QUOTE PROOF:
% "Let \(\delta(M)\) denote the depth of \(M\) and let \(i(M)\) denote
%  the smallest integer \(i\) such that \(\Ext^i_R(R/\mathfrak m, M)\)
%  is nonzero. We will see in a moment that \(i(M) < \infty\).
%  By Lemma REF we have
%  \(\delta(M) = 0\) if and only if \(i(M) = 0\), because
%  \(\mathfrak m \in \text{Ass}(M)\) exactly means
%  that \(i(M) = 0\). Hence if \(\delta(M)\) or \(i(M)\) is \(> 0\), then we may
%  choose \(x \in \mathfrak m\) such that (a) \(x\) is a nonzerodivisor
%  on \(M\), and (b) \(\text{depth}(M/xM) = \delta(M) - 1\).
%  Consider the long exact sequence
%  of Ext-groups associated to the short exact sequence
%  \(0 \to M \to M \to M/xM \to 0\) by Lemma REF:
%  [diagram of Ext long exact sequence]
%  Since \(x \in \mathfrak m\) all the maps $\Ext^i_R(\kappa, M)
%  \to \Ext^i_R(\kappa, M)$ are zero, see
%  Lemma REF.
%  Thus it is clear that \(i(M/xM) = i(M) - 1\). Induction on
%  \(\delta(M)\) finishes the proof."
\begin{proof}
  
  
\leanok
  Write \(\delta := \text{depth}(M)\) and \(i := \inf\{ j : \text{Ext}^j_R(\kappa, M) \neq 0 \}\).
  The base case \(\delta = 0\) is exactly \(\text{Hom}_R(\kappa, M) \neq 0\) (some element of
  \(M\) is annihilated by \(\mathfrak{m}\), i.e.\ \(\mathfrak{m} \in \text{Ass}(M)\), i.e.\
  \(\mathfrak{m}\) contains a zero-divisor on \(M\)), which is \(i = 0\). For the inductive step
  pick \(x \in \mathfrak{m}\) a non-zero-divisor on \(M\) with
  \(\text{depth}(M/xM) = \delta - 1\). The short exact sequence
  \(0 \to M \xrightarrow{x} M \to M/xM \to 0\) induces a long exact sequence of \(\text{Ext}^*_R(\kappa,
  -)\). Multiplication by \(x\) annihilates each \(\text{Ext}^j_R(\kappa, M)\) (because
  \(x \in \mathfrak{m}\) and \(\kappa\) is annihilated by \(\mathfrak{m}\)), so the long exact
  sequence breaks into short pieces giving
  \(\text{Ext}^{j-1}_R(\kappa, M/xM) \cong \text{Ext}^j_R(\kappa, M)\) at the relevant
  indices, and the smallest non-vanishing index drops by \(1\). By induction
  \(i(M/xM) = i - 1 = \delta - 1\), so \(i = \delta\).
\end{proof}

\section{Projective dimension}
\label{sec:ab_pd}


The two basic facts about \(\text{pd}_R(M)\) needed downstream are:
\begin{itemize}
  \item \(\text{pd}_R(M) = 0\) iff \(M\) is projective (in the local Noetherian setting,
    iff \(M\) is free);
  \item if \(0 \to K \to P \to M \to 0\) is a short exact sequence with \(P\) projective, then
    \(\text{pd}_R(M) = \text{pd}_R(K) + 1\) provided \(K \neq 0\).
\end{itemize}
Both are standard homological algebra and are recorded under
\texttt{Module.projectiveDimension\_succ\_iff} or similar in Mathlib.

\section{Depth on a short exact sequence}
\label{sec:ab_depth_ses}


% SOURCE QUOTE PROOF:
% "Use the characterization of depth using the Ext groups
%  \(\Ext^i(\kappa, N)\), see Lemma REF,
%  and use the long exact cohomology sequence
%  [diagram of Ext long exact sequence]
%  from Lemma REF."
\begin{proof}
  
  
\leanok
  \uses{lem:depth_via_ext, lem:ext_vanish_of_natCast_lt_depth, lem:natCast_add_one_le_of_le_sub_one}
  By \texttt{lem:depth_via_ext}, \(\text{depth}(N')\) (resp.\ \(\text{depth}(N)\),
  \(\text{depth}(N'')\)) is the smallest index at which \(\text{Ext}^*_R(\kappa, -)\) is
  nonzero on \(N'\) (resp.\ \(N\), \(N''\)). Apply the long exact \(\text{Ext}^*_R(\kappa, -)\)
  sequence to \(0 \to N' \to N \to N'' \to 0\):
  \[
    \cdots \to \text{Ext}^{i-1}_R(\kappa, N'')
    \to \text{Ext}^i_R(\kappa, N')
    \to \text{Ext}^i_R(\kappa, N)
    \to \text{Ext}^i_R(\kappa, N'')
    \to \text{Ext}^{i+1}_R(\kappa, N') \to \cdots
  \]
  Each of the three inequalities is then a direct read-off of the fact that if two
  of the three modules in a length-3 chunk vanish then the middle vanishes (for (1)) /
  the boundary terms vanish (for (2),(3)). Compare with the diagonal vanishing argument in
  the proof of \texttt{lem:depth_via_ext}.
\end{proof}

\section{The Auslander--Buchsbaum formula}
\label{sec:ab_main}

\begin{theorem}

\leanok
  [Auslander--Buchsbaum formula]
  \label{thm:auslander_buchsbaum}
  \lean{RingTheory.auslander_buchsbaum_formula}
  
  % SOURCE: [Stacks Project], tag 090V (proposition-Auslander-Buchsbaum)
  % (read from references/stacks-algebra.tex, L27207-L27280).
  % SOURCE QUOTE:
  % "Let \(R\) be a Noetherian local ring. Let \(M\) be a nonzero finite \(R\)-module
  %  which has finite projective dimension \(\text{pd}_R(M)\). Then we have
  %  \(\)
  %  \text{depth}(R) = \text{pd}_R(M) + \text{depth}(M)
  %  \(\)"
  \textit{Source: [Stacks Project], tag 090V (proposition-Auslander-Buchsbaum); cf.\
  Matsumura, \emph{Commutative Ring Theory}, Theorem 19.1; Auslander--Buchsbaum,
  \emph{Homological dimension in local rings}, 1957.}
  Let \((R, \mathfrak{m})\) be a Noetherian local ring and let \(M\) be a nonzero finitely
  generated \(R\)-module of finite projective dimension. Then
  \[
    \text{pd}_R(M) \;+\; \text{depth}(M) \;=\; \text{depth}(R).
  \]
\end{theorem}

% SOURCE QUOTE PROOF:
% "We prove this by induction on \(\text{depth}(M)\). The most interesting
%  case is the case \(\text{depth}(M) = 0\). In this case, let
%  \(\)
%  0 \to R^{n_e} \to R^{n_{e-1}} \to \ldots \to R^{n_0} \to M \to 0
%  \(\)
%  be a minimal finite free resolution, so \(e = \text{pd}_R(M)\).
%  By Lemma REF we may assume all matrix
%  coefficients of the maps in the complex are contained in the maximal
%  ideal of \(R\). Then on the one hand, by
%  Proposition REF we see that
%  \(\text{depth}(R) \geq e\). On the other hand, breaking the long
%  exact sequence into short exact sequences
%  \begin{align*}
%  0 \to R^{n_e} \to R^{n_{e - 1}} \to K_{e - 2} \to 0,\\
%  0 \to K_{e - 2} \to R^{n_{e - 2}} \to K_{e - 3} \to 0,\\
%  \ldots,\\
%  0 \to K_0 \to R^{n_0} \to M \to 0
%  \end{align*}
%  we see, using Lemma REF, that
%  \begin{align*}
%  \text{depth}(K_{e - 2}) \geq \text{depth}(R) - 1,\\
%  \text{depth}(K_{e - 3}) \geq \text{depth}(R) - 2,\\
%  \ldots,\\
%  \text{depth}(K_0) \geq \text{depth}(R) - (e - 1),\\
%  \text{depth}(M) \geq \text{depth}(R) - e
%  \end{align*}
%  and since \(\text{depth}(M) = 0\) we conclude \(\text{depth}(R) \leq e\).
%  This finishes the proof of the case \(\text{depth}(M) = 0\).
%
%  Induction step. If \(\text{depth}(M) > 0\), then we pick \(x \in \mathfrak m\)
%  which is a nonzerodivisor on both \(M\) and \(R\). This is possible, because
%  either \(\text{pd}_R(M) > 0\) and \(\text{depth}(R) > 0\) by the aforementioned
%  Proposition REF or \(\text{pd}_R(M) = 0\) in which
%  case \(M\) is finite free hence also \(\text{depth}(R) = \text{depth}(M) > 0\).
%  Thus \(\text{depth}(R \oplus M) > 0\) by Lemma REF
%  (for example) and we can find an \(x \in \mathfrak m\) which is a nonzerodivisor
%  on both \(R\) and \(M\). Let
%  \(\)
%  0 \to R^{n_e} \to R^{n_{e-1}} \to \ldots \to R^{n_0} \to M \to 0
%  \(\)
%  be a minimal resolution as above. An application of the snake lemma
%  shows that
%  \(\)
%  0 \to (R/xR)^{n_e} \to (R/xR)^{n_{e-1}} \to \ldots \to (R/xR)^{n_0} \to
%  M/xM \to 0
%  \(\)
%  is a minimal resolution too. Thus \(\text{pd}_R(M) = \text{pd}_{R/xR}(M/xM)\).
%  By Lemma REF we have
%  \(\text{depth}(R/xR) = \text{depth}(R) - 1\) and
%  \(\text{depth}(M/xM) = \text{depth}(M) - 1\).
%  Till now depths have all been depths as \(R\) modules, but we observe that
%  \(\text{depth}_R(M/xM) = \text{depth}_{R/xR}(M/xM)\) and similarly for \(R/xR\).
%  By induction hypothesis we see that the
%  Auslander-Buchsbaum formula holds for \(M/xM\) over \(R/xR\). Since the
%  depths of both \(R/xR\) and \(M/xM\) have decreased by one and the projective
%  dimension has not changed we conclude."
\begin{proof}
  
  
\leanok
  \uses{lem:depth_short_exact_sequence, lem:auslander_buchsbaum_formula_succ_pd,
    lem:depth_eq_of_linearEquiv, lem:depth_pi_const_eq_depth_of_nonempty}
  We induct on \(\text{depth}(M)\).

  \emph{Base case \(\text{depth}(M) = 0\).} Let \(e := \text{pd}_R(M)\) and pick a minimal
  finite free resolution
  \[
    0 \longrightarrow R^{n_e} \longrightarrow R^{n_{e-1}} \longrightarrow
    \cdots \longrightarrow R^{n_0} \longrightarrow M \longrightarrow 0
  \]
  with all transition maps having matrix coefficients in \(\mathfrak{m}\) (this is the
  ``minimal-resolution-after-trimming-trivial-summands'' normalisation, cf.\ Stacks
  \texttt{lemma-add-trivial-complex}). On the one hand, the ``what is exact'' criterion
  (Stacks tag 00MF, \texttt{proposition-what-exact}, expressing exactness of a sequence of
  finite free modules in terms of depths of ideals of \(r\)-minors) gives
  \(\text{depth}(R) \geq e\). On the other hand, splitting the resolution into short exact
  sequences
  \[
    0 \to R^{n_e} \to R^{n_{e-1}} \to K_{e-2} \to 0, \quad
    0 \to K_{e-2} \to R^{n_{e-2}} \to K_{e-3} \to 0, \quad \dots, \quad
    0 \to K_0 \to R^{n_0} \to M \to 0,
  \]
  and applying \texttt{lem:depth_short_exact_sequence} iteratively,
  yields the chain of inequalities
  \(\text{depth}(K_{e-2}) \geq \text{depth}(R) - 1\),
  \(\text{depth}(K_{e-3}) \geq \text{depth}(R) - 2\), \dots,
  \(\text{depth}(K_0) \geq \text{depth}(R) - (e-1)\), and finally
  \(\text{depth}(M) \geq \text{depth}(R) - e\). Since \(\text{depth}(M) = 0\) we deduce
  \(\text{depth}(R) \leq e\). Combined with the previous inequality, $\text{depth}(R) = e =
  \text{pd}_R(M) + 0 = \text{pd}_R(M) + \text{depth}(M)$.

  \emph{Inductive step \(\text{depth}(M) > 0\).} We claim there exists \(x \in \mathfrak{m}\)
  which is a non-zero-divisor on both \(R\) and \(M\). Indeed: either \(\text{pd}_R(M) > 0\), in
  which case \(\text{depth}(R) > 0\) by the ``what is exact'' criterion applied to the start
  of a minimal resolution, or \(\text{pd}_R(M) = 0\) in which case \(M\) is finite free and
  \(\text{depth}(R) = \text{depth}(M) > 0\). Either way both \(\text{depth}(R) > 0\) and
  \(\text{depth}(M) > 0\), hence \(\text{depth}(R \oplus M) > 0\) by
  \texttt{lem:depth_short_exact_sequence}, so a common non-zero-divisor \(x \in
  \mathfrak{m}\) exists. Take a minimal finite free resolution of \(M\) as above; the snake
  lemma applied to multiplication by \(x\) gives a minimal finite free resolution
  \[
    0 \to (R/xR)^{n_e} \to (R/xR)^{n_{e-1}} \to \cdots \to (R/xR)^{n_0} \to M/xM \to 0
  \]
  of \(M/xM\) over \(R/xR\), of the same length \(e\). So
  \(\text{pd}_{R/xR}(M/xM) = \text{pd}_R(M)\). By Stacks
  \texttt{lemma-depth-drops-by-one}, both \(\text{depth}(R/xR) = \text{depth}(R) - 1\) and
  \(\text{depth}(M/xM) = \text{depth}(M) - 1\), and one verifies
  \(\text{depth}_R(M/xM) = \text{depth}_{R/xR}(M/xM)\) (regular sequences in \(\mathfrak{m}\)
  correspond to regular sequences in \(\mathfrak{m}/(x)\)). Apply the inductive hypothesis to
  \(M/xM\) over \(R/xR\):
  \[
    \text{pd}_{R/xR}(M/xM) + \text{depth}(M/xM) = \text{depth}(R/xR).
  \]
  Substituting yields
  \(\text{pd}_R(M) + (\text{depth}(M) - 1) = \text{depth}(R) - 1\), i.e.\
  \(\text{pd}_R(M) + \text{depth}(M) = \text{depth}(R)\), as required.
\end{proof}


\subsection{Inductive-step substrate: per-gap decomposition of
  \texorpdfstring{\(\text{pd}(M) = k+1\)}{pd(M) = k+1}}
\label{subsec:succ_pd_gap_sequence}

\paragraph{Status: RESOLVED iter-202.}
\texttt{lem:auslander_buchsbaum_formula_succ_pd} and the main formula
\cref{thm:auslander_buchsbaum} are both closed axiom-clean as of
iter-202, by a \(\texttt{Nat}\)-induction restructuring
(\(\texttt{induction}\,k\,\texttt{generalizing}\,M\)): the inductive
step uses the kernel SES with
\(\texttt{depth\_of\_short\_exact}\) parts (2),(3) directly (no
exactness criterion), and the base case \(\texttt{pd}\,M = 1\) uses
the iter-201 matrix-collapse helper
(\(\texttt{ext\_comp\_mk}_0\texttt{\_ofHom\_eq\_zero\_of\_entries\_mem\_annihilator}\))
together with the long exact sequence. \textbf{Gap (2) (Stacks 00MF,
the ``what is exact'' criterion) was OBVIATED} by this Path~B route ---
no Buchsbaum--Eisenbud exactness criterion was needed. \textbf{All
gaps (1)--(4) are CLOSED or OBVIATED.} The decomposition below is
retained as the historical dependency analysis that motivated the
substrate builds; it is no longer an open work-list.

\medskip
The inductive step
\texttt{lem:auslander_buchsbaum_formula_succ_pd}
of the Auslander--Buchsbaum formula was originally gated, against an
earlier pinned Mathlib snapshot, by three absent substrate ingredients
(gaps (1)--(3)) plus a fourth in hand
(gap (4), \texttt{lem:depth_drops_by_one}). This subsection records the
decomposition that the proof block of
\texttt{lem:auslander_buchsbaum_formula_succ_pd} folded into four bullets:
its purpose was to make the dependency graph and the per-gap effort
estimates explicit, so that the substrate gaps could be formalised
independently (and where the dependency graph permitted, in parallel).

\paragraph{Sequencing constraint.} The three open gaps are
\emph{not} mutually independent at the substrate level. Gap (3)
(snake-lemma carving of the minimal resolution by an
\(M\)-regular element \(x \in \mathfrak m\)) needs gap (1)
(minimal-resolution normalisation: every finite free resolution of
a finite module over a Noetherian local ring trims to a minimal one
of the same projective dimension) as input, because the snake-lemma
output is only known to be a \emph{minimal} resolution of \(M/xM\)
over \(R/(x)\) of the same length when the input resolution is itself
minimal. Gap (2) (Stacks 00MF, the ``what is exact'' criterion in
terms of depths of \(r\)-minors of the transition matrices) is
independent of gaps (1) and (3) at the substrate level: it is a
self-contained criterion on a complex of finite free modules and
feeds the closure assembly without going through gap (1) or gap (3).

\paragraph{Dependency diagram.} Read top-to-bottom as ``provides
input to'':
\begin{verbatim}
                gap (1)
       minimal-resolution carving
   (Stacks lemma-add-trivial-complex)
                  |
                  v
                gap (3)
        snake lemma on minimal
        resolution under x in m
                  |
                  v
   +--------------+--------------+
   |   L1299 closure assembly    |   <--   gap (2)
   |   (succ-pd inductive step)  |        Stacks 00MF
   +-----------------------------+        "what is exact"
                  ^
                  |
                gap (4)
       depth drops by one (CLOSED;
       see the lemma above)
\end{verbatim}
Gap (4) is already formalised
(\texttt{lem:depth_drops_by_one}, the project's
\texttt{depth\_quotSMulTop\_succ\_eq\_depth\_of\_isSMulRegular})
and feeds the closure assembly directly. Gap (2) feeds the closure
assembly independently of the gap (1) \(\to\) gap (3) chain.

\paragraph{Per-gap effort estimates.} The following table records,
for each gap, the estimated formalisation effort (lines of code in
the project's Lean source style), a coarse iter budget, the
Mathlib-at-\texttt{b80f227} status, and substrate dependencies on
other gaps:
\begin{center}
  \begin{tabular}{|c|l|c|c|c|c|}
    \hline
    Gap & Statement                              & LOC est.  & Iters
        & Mathlib status     & Depends on \\
    \hline
    (1) & Minimal-resolution carving             & 40--80    & 1--2
        & per-syzygy step CLOSED iter-199
                                                                & none \\
    (2) & Stacks 00MF (``what is exact'')        & 150--200  & 2--3
        & absent             & none \\
    (3) & Snake lemma on minimal resolution      & ---       & ---
        & \textbf{OBVIATED iter-200} & --- \\
    (4) & Depth drops by one                     & ---       & ---
        & \textbf{CLOSED}    & --- \\
    \hline
  \end{tabular}
\end{center}
The gap (1) estimate is anchored to Bruns--Herzog,
\emph{Cohen--Macaulay Rings}, Construction~1.5.18 together with the
minimal-resolution induction (\S1.5): the trimming algorithm is a
finite-rank linear-algebra computation followed by induction on the
length of the resolution. The gap (2) estimate reflects the chain
Bruns--Herzog Theorem~1.4.13 \(\to\) Eagon--Northcott
\(\to\) Buchsbaum--Eisenbud's exactness criterion (Stacks 00MF /
\texttt{proposition-what-exact}), which characterises exactness of a
complex of finite free modules in terms of depths of the ideals of
\(r\)-minors of the transition matrices. Gap (2) is the largest of
the three open substrates and is a natural candidate for an upstream
Mathlib contribution rather than a project-internal proof.

\paragraph{Gap (3) OBVIATED iter-200.} The iter-200 ALIGN\_WITH\_MATHLIB
pivot (\ref{subsec:ab_gap1_haspdlt_pivot}) replaces the literal
\(\texttt{ChainComplex}\,\mathbb N\) carrier with the abstract
\(\texttt{HasProjectiveDimensionLT}\) SES-descent pattern; under this
pattern, minimality of the resolution is \emph{not} used as input
anywhere in the inductive step. Consequently the snake-lemma carving
of the minimal resolution (gap (3)) is no longer a substrate
requirement: only gap (2) (\(\texttt{pd}\,M > 0 \Rightarrow
\texttt{depth}\,M < \texttt{depth}\,R\)) remains the binding gap for
the closure of \texttt{lem:auslander_buchsbaum_formula_succ_pd}, and the
per-syzygy step of gap (1) is consumed only through the SES at
\(0 \to \ker f \to R^n \to M \to 0\) (no minimality input required).
See \ref{subsec:ab_gap1_haspdlt_pivot} for the SES-descent details.

\paragraph{Gap (2) Stacks 00MF proof recipe.} The
Buchsbaum--Eisenbud exactness criterion (Stacks 00MF /
\texttt{proposition-what-exact}) reads: a complex
\(\dots \to R^{r_2} \xrightarrow{\varphi_2} R^{r_1}
 \xrightarrow{\varphi_1} R^{r_0}\) of finite free modules over a
Noetherian local ring \(R\) is exact at \(R^{r_i}\) (for \(i \geq 1\))
if and only if \(\texttt{rank}\,\varphi_{i+1} + \texttt{rank}\,\varphi_i
= r_i\) and \(\texttt{depth}\,I(\varphi_i) \geq i\), where
\(I(\varphi)\) denotes the ideal of \(r\)-minors of \(\varphi\) for
\(r = \texttt{rank}\,\varphi\). The recipe for an iter-201+ prover:
(a) define \(I(\varphi) \subseteq R\) as the ideal of all
\(\texttt{rank}\,\varphi\)-minors of \(\varphi\) (Mathlib supports
this via \(\texttt{Matrix.minors}\) and \(\texttt{Ideal.span}\) over the
finite index set of \(\binom{n}{r}\) tuples); (b) prove the depth
lower bound via induction on \(i\) using a Koszul-complex argument
on the regular sequence implicit in the exactness data; (c) derive
\(\texttt{pd}\,M > 0 \Rightarrow \texttt{depth}\,M < \texttt{depth}\,R\)
as a corollary by taking the start \(\dots \to R^{r_1} \to R^{r_0} \to
M\) of a minimal resolution and applying the criterion at \(i = 1\).
Estimated \(\sim 150-200\) LOC project-side; Mathlib upstream PR
candidate.

\paragraph{Closure assembly (iter-200 ALIGN\_WITH\_MATHLIB pivot).} The
iter-200 substrate landing in \ref{subsec:ab_gap1_haspdlt_pivot}
replaces the original (1)+(3)-dependent assembly with an SES-descent
recipe: given an \(M\) with \(\texttt{pd}_R(M) = k+1\), the
per-syzygy step of gap (1)
(\texttt{lem:exists_minimalSurjection_finite_localRing}) supplies a
surjection \(f : R^n \twoheadrightarrow M\); the SES
\(0 \to \ker f \to R^n \to M \to 0\) plus the iter-200 helpers
\(\texttt{hasProjectiveDimensionLT\_ker\_of\_surjection}\) +
\(\texttt{hasProjectiveDimensionLT\_succ\_of\_hasProjectiveDimensionLT\_ker}\)
give \(\texttt{pd}_R(\ker f) = k\) exactly; gap (2) (the
Buchsbaum--Eisenbud criterion at the start of a minimal resolution)
delivers \(\texttt{depth}\,M < \texttt{depth}\,R\); the iter-200
helper
\(\texttt{depth\_ker\_ge\_min\_of\_surjection\_finite\_localRing}\)
gives \(\texttt{depth}\,\ker f \geq \min(\texttt{depth}\,R,
\texttt{depth}\,M + 1) = \texttt{depth}\,M + 1\); a complementary LES
upper bound (or the IH applied to \(\ker f\) under the strict
descent of \(\texttt{pd}\)) gives the matching upper bound. Combine
to get \(\texttt{depth}\,\ker f = \texttt{depth}\,M + 1\), then apply
the IH (\(\texttt{pd}\,\ker f = k\)) for
\(k + \texttt{depth}\,\ker f = \texttt{depth}\,R\), which rearranges
to \((k+1) + \texttt{depth}\,M = \texttt{depth}\,R\) as required.
Estimated closure-assembly cost (gaps in hand): \(\sim 50-80\) LOC
inductive code.

\paragraph{Iter budget refinement.} Aggregating the per-gap
estimates with the closure assembly: gap (1) at 1--2 iters, gap (2)
at 2--3 iters, gap (3) at 1--2 iters, and the closure assembly
at 1 iter gives a total remaining iter budget of \textbf{5--8 iters}
for \texttt{lem:auslander_buchsbaum_formula_succ_pd}. The schedule may
parallelise across the gap (2) branch and the gap (1) \(\to\) gap (3)
chain --- so wall-clock could be as low as
\(\max(\text{gap (2)},\, \text{gap (1)} + \text{gap (3)}) +
\text{closure}\), i.e.\ \(\max(2\text{--}3,\, 2\text{--}4) + 1\),
yielding 3--5 iters in the parallel-best case. The substrate
sequencing constraint (gap (3) requires gap (1)) is hard; the gap
(2) branch is genuinely independent of that chain and could be
dispatched as a parallel prover lane the moment a free slot opens.
This estimate refines the earlier ``~6--12 iters'' coarse budget
recorded in \texttt{STRATEGY.md}: the refined window is somewhat
tighter on the upper end (5--8 vs.\ 6--12 sequential) and is
substantially tighter under parallel dispatch (3--5 parallel-best
vs.\ 6--12).

\subsection{Gap (1) first-step substrate (iter-199): per-syzygy
  minimal surjection from a finite free module}
\label{subsec:ab_gap1_first_step}


This lemma is the \emph{per-syzygy} step of gap (1)'s
minimal-resolution carving: given a finite \(R\)-module \(M\)
over a local ring \(R\), it produces a minimal surjection
\(R^n \twoheadrightarrow M\). Iter-200's ALIGN\_WITH\_MATHLIB pivot
(\ref{subsec:ab_gap1_haspdlt_pivot}) replaces the original
\(\mathbb N\)-recursive \(\texttt{ChainComplex}\) assembly target
with abstract \(\texttt{HasProjectiveDimensionLT}\) SES descent: only
the per-syzygy SES is consumed; the iterated chain-complex carrier
is not needed.

\subsection{Iter-200 ALIGN\_WITH\_MATHLIB pivot: \(\texttt{HasProjectiveDimensionLT}\) SES descent}
\label{subsec:ab_gap1_haspdlt_pivot}

The original iter-199 plan for closing
\texttt{lem:auslander_buchsbaum_formula_succ_pd} carried the
\(\mathbb N\)-recursive minimal free resolution as a literal
\(\texttt{ChainComplex}\,\mathbb N\,(\texttt{ModuleCat}\,R)\) object,
proved finite-rank-free at each level, and then attacked the
inductive step via the snake lemma (gap (3)). The
\texttt{ab-natrecursive} mathlib-analogist verdict (iter-200,
persistent file \texttt{analogies/ab-natrecursive.md}) rejected this
route as 3--4\(\times\) over budget and gated on a Mathlib-absent
termination-of-syzygy-tower lemma. The recommended alignment is to
descend pd via Mathlib's
\(\texttt{CategoryTheory.HasProjectiveDimensionLT}\) predicate, which
takes \(\texttt{ShortComplex.ShortExact}\) input and an extremal
projectivity hypothesis on \(X_1\) or \(X_3\) to give
\(\texttt{HasProjectiveDimensionLT}\) on the other piece. No
\(\mathbb N\)-recursive carrier and no minimality content are
required.

\paragraph{Sequencing for closure.} The four iter-200 helpers above
plus gap (2) (Stacks 00MF, the \(\texttt{depth}\,M < \texttt{depth}\,R\)
input from \texttt{pd}-positivity) suffice for the closure of
\texttt{lem:auslander_buchsbaum_formula_succ_pd}: the per-syzygy step
(\texttt{lem:exists_minimalSurjection_finite_localRing}) supplies \(f\);
\texttt{lem:hasProjectiveDimensionLT_succ_of_projectiveDimension_eq}
+ \texttt{lem:hasProjectiveDimensionLT_ker_of_surjection} +
\texttt{lem:hasProjectiveDimensionLT_succ_of_hasProjectiveDimensionLT_ker}
give \(\texttt{pd}\,\ker f = k\) exactly;
\texttt{lem:depth_ker_ge_min_of_surjection_finite_localRing} plus gap (2)
plus a complementary LES upper bound (or the IH) give
\(\texttt{depth}\,\ker f = \texttt{depth}\,M + 1\); the inductive
hypothesis on \(\ker f\) (at pd \(=k\)) plus the rearrangement of
\((k + (\texttt{depth}\,M+1)) = \texttt{depth}\,R\) closes. The
binding remaining substrate is gap (2) (\(\sim 150\)--\(200\) LOC,
Mathlib upstream PR candidate per the recipe in
\ref{subsec:succ_pd_gap_sequence}).

\paragraph{Iter-201 Path B closure recipe (LES-injectivity +
matrix-collapse).} The iter-201 \texttt{ab-stacks00mf} mathlib-
analogist verdict identifies an alternative closure that sidesteps
gap (2): instead of building Stacks 00MF as a standalone substrate,
split \texttt{lem:auslander_buchsbaum_formula_succ_pd} into an
\emph{inductive step} (\(\texttt{pd}\,M = k+1, k \geq 1\)) and a
\emph{base case} (\(\texttt{pd}\,M = 1\)). Each piece uses only
already-existing LES infrastructure plus one new
\(\sim 50\)--\(80\) LOC matrix-collapse helper.

\textit{Inductive step \(\texttt{pd}\,M = k+1, k \geq 1\)
\((\sim 50\)--\(80\) LOC).} From
\texttt{lem:hasProjectiveDimensionLT_ker_of_surjection} +
the contrapositive of
\texttt{lem:hasProjectiveDimensionLT_succ_of_hasProjectiveDimensionLT_ker},
\(\texttt{pd}\,\ker f = k\) exactly. Strong induction gives
\(k + \texttt{depth}\,\ker f = \texttt{depth}\,R\), hence
\(\texttt{depth}\,\ker f < \texttt{depth}\,R\) since \(k \geq 1\).
Lower bound \(\texttt{depth}\,M \geq \texttt{depth}\,\ker f - 1\):
\(\texttt{depth\_of\_short\_exact}\) part (2). Upper bound
\(\texttt{depth}\,M \leq \texttt{depth}\,\ker f - 1\): by
contradiction --- assume \(\texttt{depth}\,M \geq \texttt{depth}\,\ker f\),
look at the LES of \(\texttt{Ext}^*(\kappa, -)\) at
\(i = \texttt{depth}\,\ker f\), infer
\(\texttt{Ext}^{\texttt{depth}\,\ker f}(\kappa, R^n) \neq 0\)
(forcing \(\texttt{depth}\,R \leq \texttt{depth}\,\ker f\)),
contradicting the IH.

\textit{Base case \(\texttt{pd}\,M = 1\) \((\sim 80\)--\(120\) LOC).}
\(\ker f\) is projective (\(\texttt{pd}\,\ker f = 0\)), hence free
over the local Noetherian \(R\) by Mathlib's
\(\texttt{Module.free\_of\_flat\_of\_isLocalRing}\) +
\(\texttt{Module.Flat.of\_projective}\). So \(\ker f \cong R^m\) with
\(m \geq 1\). The inclusion \(\ker f \hookrightarrow R^n\) is an
\((n \times m)\)-matrix \(A\) with entries in
\(\texttt{Ann}(\kappa) = \mathfrak m\) (from
\texttt{lem:exists_minimalSurjection_finite_localRing}'s minimality
clause \(\ker f \subseteq \mathfrak m \cdot R^n\)).

The \textbf{matrix-collapse helper}
(CLOSED iter-201, 4 axiom-clean private decls at
\texttt{AuslanderBuchsbaum.lean} L1418--L1517):
\[
  \forall p \geq 0,\quad \forall e \in \texttt{Ext}^p(\kappa, R^n),\quad
  e \circ \texttt{Ext.mk}_0\,(\texttt{ofHom}\,A) = 0
  \quad \text{in } \texttt{Ext}^{p+0}(\kappa, R^n)
  \text{ when } A_{ij} \in \texttt{Ann}_R\,\kappa.
\]
The 4 helpers are \(\texttt{Module.elemMap}\) (the elementary map
sending \(\texttt{Pi.single}\,j\,1 \mapsto \texttt{Pi.single}\,i\,1\)
and other basis vectors to zero), \(\texttt{Module.elemMap\_apply}\),
\(\texttt{Module.linearMap\_finFunR\_matrix\_decomp}\)
(\(A = \sum_{i,j} A(\texttt{Pi.single}\,j\,1)\,i \cdot
   \texttt{elemMap}_{ij}\)), and the closure helper
\(\texttt{Module.ext\_comp\_mk}_0\texttt{\_ofHom\_eq\_zero\_of\_entries\_mem\_annihilator}\)
itself. The realised API path uses Mathlib's
\(\texttt{ModuleCat.hom\_sum}\) + \(\texttt{Ext.mk}_0\texttt{\_sum}\)
+ \(\texttt{Ext.comp\_sum}\) and the scalar siblings
\(\texttt{ModuleCat.hom\_smul}\) + \(\texttt{Ext.mk}_0\texttt{\_smul}\)
+ \(\texttt{Ext.comp\_smul}\) (\emph{not}
\(\texttt{Ext.bilinearCompOfLinear}\) as originally projected; the
mathematical content is identical). Each summand reduces to a scalar-
bilinear form killed by the existing
\(\texttt{ext\_smul\_eq\_zero\_of\_mem\_annihilator}\) (L229).

With matrix-collapse, the LES at \(i = \texttt{depth}\,R - 1\)
gives
\[
  \texttt{Ext}^{\texttt{depth}\,R - 1}(\kappa, M)
  \xrightarrow{\;\delta\;}
  \texttt{Ext}^{\texttt{depth}\,R}(\kappa, \ker f)
  \xrightarrow{\;(- \circ A)\;}
  \texttt{Ext}^{\texttt{depth}\,R}(\kappa, R^n)
\]
with the second arrow zero. If
\(\texttt{depth}\,M \geq \texttt{depth}\,R\), the source is zero,
hence \(\texttt{Ext}^{\texttt{depth}\,R}(\kappa, \ker f) = 0\) ---
but \(\texttt{Ext}^{\texttt{depth}\,R}(\kappa, \ker f) \simeq
\texttt{Ext}^{\texttt{depth}\,R}(\kappa, R^m) \neq 0\)
(\(m \geq 1\), \(\texttt{depth}\,R^m = \texttt{depth}\,R\)).
Contradiction. So \(\texttt{depth}\,M < \texttt{depth}\,R\), and the
base case
\(\texttt{pd}\,M + \texttt{depth}\,M = \texttt{depth}\,R\) follows
from \(\texttt{depth}\,\ker f = \texttt{depth}\,R\) (since \(\ker f\)
free).

\textbf{Total Path B cost: \(\sim 130\)--\(200\) LOC} comparable to
Path A but with better composition --- every new lemma reuses
existing axiom-clean primitives, and the matrix-collapse helper is
independently useful for any future ``minimal matrix annihilates
\(\texttt{Ext}^*(\kappa, -)\)'' argument. Persistent file:
\texttt{analogies/ab-stacks00mf.md} (iter-201).

\paragraph{Iter-201 status update: matrix-collapse substrate landed;
closure body deferred to iter-202.} The 4 axiom-clean matrix-collapse
helpers are in place (L1418--L1517, all private under
\(\texttt{RingTheory.Module}\)). What remains (iter-202 dispatch) is
the body of \texttt{lem:auslander_buchsbaum_formula_succ_pd} itself:
restructure as \(\texttt{induction } k \texttt{ generalizing } M\),
assemble the base case from matrix-collapse + LES at
\(i = \texttt{depth}\,R - 1\) (\(\sim 50\)--\(80\) LOC of LES
bookkeeping + linear-equivalence transport of matrix-collapse to the
subtype inclusion \(\ker f \hookrightarrow R^n\) through the basis
iso \(R^m \simeq \ker f\)), and assemble the inductive step
\(k \geq 1\) from
\(\texttt{depth\_of\_short\_exact}\) parts (2) and (3) directly ---
\emph{no matrix-collapse needed for the inductive step}, per the
iter-201 prover's discovery (\(\texttt{depth\_of\_short\_exact}\,(3)\)
gives \(\texttt{depth}\,M + 1 \leq \texttt{depth}\,\ker f\) which
combined with IH \(k + \texttt{depth}\,\ker f = \texttt{depth}\,R\)
closes the inductive step in \(\sim 30\)--\(50\) LOC). Total iter-202
closure cost: \(\sim 80\)--\(120\) LOC.

The iter-202 Lane AB dispatch also discharges the iter-201 plan
agent's option (1) commitment: remove \texttt{private} from
\texttt{lem:auslander_buchsbaum_formula_succ_pd} once the body closes,
and additionally promote \(\texttt{isDomain\_of\_regularLocal}\)
(currently private at L2657) and
\(\texttt{regularLocal\_quotient\_isRegularLocal\_of\_notMemSq}\)
(currently private at L2293) to public so that
\texttt{CodimOneExtension.lean} can consume them cross-file as the
project-local Stacks 00NQ witness and the
\(A / (f_1) \to \texttt{IsRegularLocalRing}\) preservation helper for
Lane COE Step A1 (see \(\cref{subsec:stage6_iib_substrate_iter200}\)
on the COE side).

\section{Substrate helper declarations (1-to-1 Lean coverage)}
\label{sec:ab_substrate_helpers}

This section records the internal helper declarations on which the public
results above rest. Each is a self-contained, sorry-free step of the Lean
development; the prose below states what each helper asserts and how it feeds
the corresponding public result.

\subsection{Regular-sequence and domain substrate for regular local rings}
\label{subsec:ab_regular_domain_substrate}

\begin{lemma}\leanok
[Regular sequences are bounded by the Krull dimension]
  \label{lem:length_le_ringKrullDim_of_isRegular}
  \lean{RingTheory.CohenMacaulay.length_le_ringKrullDim_of_isRegular}
  For a Noetherian local ring \(R\), every \(R\)-regular sequence \(rs\) has length at
  most \(\dim(R)\). This is the upper-bound half of \cref{cor:regular_cohen_macaulay},
  obtained from the dimension identity \(\dim(R / (rs)) + |rs| = \dim(R)\) together with
  \(\dim(R / (rs)) \geq 0\).
\end{lemma}
\begin{proof}

\leanok
  Proved directly in Lean.
\end{proof}

\begin{lemma}\leanok
[Cotangent image of \(x \in \mathfrak m \setminus \mathfrak m^2\)]
  \label{lem:toCotangent_ne_zero_of_not_mem_sq}
  \lean{RingTheory.CohenMacaulay.toCotangent_ne_zero_of_not_mem_sq}
  For a local ring \((R, \mathfrak m)\) and \(x \in \mathfrak m\) with \(x \notin \mathfrak m^2\),
  the image of \(x\) in the cotangent space \(\mathfrak m / \mathfrak m^2\) is nonzero. This is
  the positivity input for the cotangent dimension-drop.
\end{lemma}
\begin{proof}

\leanok
  Proved directly in Lean, from the criterion \(x \mapsto 0\) in the cotangent space
  iff \(x \in \mathfrak m^2\).
\end{proof}

\begin{lemma}\leanok
[Cotangent dimension drops by one modulo \(x\)]
  \label{lem:finrank_cotangentSpace_quot_span_singleton_succ}
  \lean{RingTheory.CohenMacaulay.finrank_cotangentSpace_quot_span_singleton_succ}
  \uses{lem:toCotangent_ne_zero_of_not_mem_sq}
  For a Noetherian local ring \((R, \mathfrak m)\) and \(x \in \mathfrak m \setminus \mathfrak m^2\),
  the cotangent space of \(R / (x)\) has dimension one less than that of \(R\):
  \[
    \dim_{\kappa'}\!\bigl(\text{CotangentSpace}(R/(x))\bigr) + 1
      = \dim_{\kappa}\!\bigl(\text{CotangentSpace}(R)\bigr),
  \]
  where \(\kappa, \kappa'\) are the residue fields of \(R\) and \(R/(x)\). The kernel of
  the induced surjection on cotangent spaces is the line spanned by the image of \(x\),
  which is one-dimensional by \cref{lem:toCotangent_ne_zero_of_not_mem_sq}.
\end{lemma}
\begin{proof}

\leanok
  Proved directly in Lean.
\end{proof}

\begin{lemma}\leanok
[Nakayama witness for a positive cotangent dimension]
  \label{lem:exists_notMemSq_of_spanFinrank_pos}
  \lean{RingTheory.CohenMacaulay.exists_notMemSq_of_spanFinrank_pos}
  For a Noetherian local ring \((R, \mathfrak m)\) with positive minimal-generator count
  of \(\mathfrak m\) (i.e.\ \(\text{spanFinrank}(\mathfrak m) \geq 1\)), there exists
  \(x \in \mathfrak m\) with \(x \notin \mathfrak m^2\). By Nakayama, \(\mathfrak m \subseteq \mathfrak m^2\)
  would force \(\mathfrak m = 0\), contradicting positivity.
\end{lemma}
\begin{proof}

\leanok
  Proved directly in Lean.
\end{proof}

\begin{lemma}\leanok
[Zero cotangent dimension forces a field]
  \label{lem:isDomain_of_isLocalRing_of_spanFinrank_maximalIdeal_eq_zero}
  \lean{RingTheory.CohenMacaulay.isDomain_of_isLocalRing_of_spanFinrank_maximalIdeal_eq_zero}
  For a Noetherian local ring \(R\) with \(\text{spanFinrank}(\mathfrak m) = 0\), the maximal
  ideal collapses to \(0\), so \(R\) is a field and hence an integral domain. This is the
  base case of \cref{lem:isDomain_of_regularLocal}.
\end{lemma}
\begin{proof}

\leanok
  Proved directly in Lean.
\end{proof}

\begin{lemma}\leanok
[Regularity descends to \(R/(x)\)]
  \label{lem:regularLocal_quotient_isRegularLocal_of_notMemSq}
  \lean{RingTheory.CohenMacaulay.regularLocal_quotient_isRegularLocal_of_notMemSq}
  \uses{lem:finrank_cotangentSpace_quot_span_singleton_succ}
  For a regular Noetherian local ring \(R\) with \(\text{spanFinrank}(\mathfrak m) = k + 1\)
  and \(x \in \mathfrak m \setminus \mathfrak m^2\), the quotient \(R / (x)\) is again a regular
  local ring, with \(\text{spanFinrank}(\mathfrak m') = k\). The cotangent dimension drop
  (\cref{lem:finrank_cotangentSpace_quot_span_singleton_succ}) plus a Krull-height bound
  give the matching dimension equality, so \(R/(x)\) is regular. Notably this avoids
  assuming \(x\) is a non-zero-divisor, so it does not depend on
  \cref{lem:isDomain_of_regularLocal}.
\end{lemma}
\begin{proof}

\leanok
  Proved directly in Lean.
\end{proof}

\begin{lemma}\leanok
[Members of a minimal prime are zero-divisors]
  \label{lem:exists_ne_zero_mul_eq_zero_of_mem_minimalPrimes}
  \lean{RingTheory.CohenMacaulay.exists_ne_zero_mul_eq_zero_of_mem_minimalPrimes}
  For a commutative ring \(R\) and a minimal prime \(\mathfrak p\) of \(R\), every
  \(x \in \mathfrak p\) is a zero-divisor: there exists \(y \neq 0\) with \(x y = 0\). This
  follows from the disjointness of a minimal prime from the non-zero-divisors.
\end{lemma}
\begin{proof}

\leanok
  Proved directly in Lean.
\end{proof}

\begin{lemma}\leanok
[\((x)\) is not a minimal prime in the inductive step]
  \label{lem:notMem_minimalPrimes_of_regularLocal_succ}
  \lean{RingTheory.CohenMacaulay.notMem_minimalPrimes_of_regularLocal_succ}
  \uses{lem:exists_ne_zero_mul_eq_zero_of_mem_minimalPrimes, lem:regularLocal_quotient_isRegularLocal_of_notMemSq}
  Let \(R\) be a regular Noetherian local ring with \(\text{spanFinrank}(\mathfrak m) = k + 1\)
  and \(x \in \mathfrak m \setminus \mathfrak m^2\), and suppose every \(k\)-dimensional regular
  local ring is a domain. Then \(\text{span}\{x\}\) is \emph{not} a minimal prime of \(R\).
  Combined with the regular descent
  (\cref{lem:regularLocal_quotient_isRegularLocal_of_notMemSq}), this is the residual
  case of the regular-local domain argument.
\end{lemma}
\begin{proof}

\leanok
  Proved directly in Lean.
\end{proof}

\begin{lemma}\leanok
[Regular local rings are domains]
  \label{lem:isDomain_of_regularLocal}
  \lean{RingTheory.CohenMacaulay.isDomain_of_regularLocal}
  \uses{lem:isDomain_of_isLocalRing_of_spanFinrank_maximalIdeal_eq_zero, lem:exists_notMemSq_of_spanFinrank_pos, lem:regularLocal_quotient_isRegularLocal_of_notMemSq, lem:notMem_minimalPrimes_of_regularLocal_succ}
  Every regular Noetherian local ring \(R\) is an integral domain. The proof is a strong
  induction on \(\text{spanFinrank}(\mathfrak m)\): the base case is the field case
  (\cref{lem:isDomain_of_isLocalRing_of_spanFinrank_maximalIdeal_eq_zero}); the inductive
  step picks \(x \in \mathfrak m \setminus \mathfrak m^2\)
  (\cref{lem:exists_notMemSq_of_spanFinrank_pos}), uses that \(R/(x)\) is regular local of
  smaller dimension (\cref{lem:regularLocal_quotient_isRegularLocal_of_notMemSq}) and a
  domain by induction, and rules out the obstruction case via
  \cref{lem:notMem_minimalPrimes_of_regularLocal_succ}.
\end{lemma}
\begin{proof}

\leanok
  Proved directly in Lean.
\end{proof}

\begin{lemma}\leanok
[A regular element outside \(\mathfrak m^2\)]
  \label{lem:exists_isSMulRegular_notMemSq_of_regularLocal_succ}
  \lean{RingTheory.CohenMacaulay.exists_isSMulRegular_notMemSq_of_regularLocal_succ}
  \uses{lem:exists_notMemSq_of_spanFinrank_pos, lem:isDomain_of_regularLocal}
  For a regular local ring \(R\) of dimension \(k + 1\), there exists
  \(x \in \mathfrak m \setminus \mathfrak m^2\) that is moreover an \(R\)-regular element. The
  Nakayama witness (\cref{lem:exists_notMemSq_of_spanFinrank_pos}) is nonzero, and in a
  domain (\cref{lem:isDomain_of_regularLocal}) every nonzero element is a
  non-zero-divisor.
\end{lemma}
\begin{proof}

\leanok
  Proved directly in Lean.
\end{proof}

\begin{lemma}\leanok
[Regular element with regular quotient]
  \label{lem:exists_isSMulRegular_quotient_isRegularLocal_succ}
  \lean{RingTheory.CohenMacaulay.exists_isSMulRegular_quotient_isRegularLocal_succ}
  \uses{lem:exists_isSMulRegular_notMemSq_of_regularLocal_succ, lem:finrank_cotangentSpace_quot_span_singleton_succ}
  For a regular local ring \(R\) of dimension \(k + 1\), there exists \(x \in \mathfrak m\)
  that is \(R\)-regular and such that \(R / (x)\) is again a regular local ring of
  dimension \(k\) (its maximal ideal has \(\text{spanFinrank} = k\)). This consolidates the
  ``pick an \(R\)-regular system-of-parameters element and drop dimension'' step.
\end{lemma}
\begin{proof}

\leanok
  Proved directly in Lean.
\end{proof}

\begin{lemma}\leanok
[Inductive step for the regular sequence]
  \label{lem:regularLocal_inductive_step}
  \lean{RingTheory.CohenMacaulay.regularLocal_inductive_step}
  \uses{lem:exists_isSMulRegular_quotient_isRegularLocal_succ}
  Let \(R\) be a regular local ring of dimension \(k + 1\), and suppose that every regular
  local ring of dimension \(k\) admits a length-\(k\) regular sequence inside its maximal
  ideal. Then \(R\) admits a length-\((k+1)\) regular sequence inside \(\mathfrak m\): extract
  the regular element \(x\) with regular quotient
  (\cref{lem:exists_isSMulRegular_quotient_isRegularLocal_succ}), lift a length-\(k\)
  regular sequence of \(R/(x)\) back to \(R\), and prepend \(x\).
\end{lemma}
\begin{proof}

\leanok
  Proved directly in Lean.
\end{proof}

\begin{lemma}\leanok
[Regular local rings have a full regular sequence]
  \label{lem:exists_isRegular_of_regularLocal}
  \lean{RingTheory.CohenMacaulay.exists_isRegular_of_regularLocal}
  \uses{lem:regularLocal_inductive_step}
  For a regular Noetherian local ring \(R\), there is an \(R\)-regular sequence inside
  \(\mathfrak m\) whose length equals \(\text{spanFinrank}(\mathfrak m) = \dim(R)\). The proof is a
  strong induction on the dimension: the base case is the empty sequence, and the
  inductive step is \cref{lem:regularLocal_inductive_step}. This is the lower-bound half
  of \cref{cor:regular_cohen_macaulay}.
\end{lemma}
\begin{proof}

\leanok
  Proved directly in Lean.
\end{proof}

\section{Corollary: regular local rings are Cohen--Macaulay}
\label{sec:ab_regular_cm}

\begin{definition}

\leanok
  [Cohen--Macaulay local ring]
  \label{def:cohen_macaulay_local}
  \lean{RingTheory.CohenMacaulay}
  \uses{def:depth}
  % SOURCE: [Stacks Project], tag 00N4 (definition-local-ring-CM)
  % (read from references/stacks-algebra.tex, L25288-L25299).
  % SOURCE QUOTE:
  % "A Noetherian local ring \(R\) is called {\it Cohen-Macaulay}
  %  if it is Cohen-Macaulay as a module over itself."
  \textit{Source: [Stacks Project], tag 00N4 (definition-local-ring-CM).}
  A Noetherian local ring \((R, \mathfrak{m})\) is called \emph{Cohen--Macaulay} if
  \(\text{depth}(R) = \dim(R)\) (equivalently, if \(R\) is Cohen--Macaulay as a module
  over itself, where Cohen--Macaulay for a finite \(R\)-module \(M\) means
  \(\dim(\text{Supp}(M)) = \text{depth}(M)\), cf.\ Stacks tag 00N3).
\end{definition}

% NOTE (iter-185 review, supersedes iter-184): the Lean assembly
% `CohenMacaulay.of_regular` calls a non-private helper
% `RingTheory.CohenMacaulay.exists_isRegular_of_regularLocal` (AuslanderBuchsbaum.lean
% L1096 — line drifted from iter-184's L944 reference) that now carries a
% structural induction scaffold via `regularLocal_inductive_step`, with the
% load-bearing sorry pushed into two named helpers: (i)
% `exists_isSMulRegular_quotient_isRegularLocal_succ` consolidating the
% Stacks 00NQ + 00NU substrate gaps in one signature, and (ii) an inline
% technical-bridge sorry inside `regularLocal_inductive_step` for the
% `R⧸(x)`-vs-`QuotSMulTop x R` linear-equiv upgrade (~10-20 LOC iter-186 close
% via `Submodule.ideal_span_singleton_smul` + `QuotSMulTop.mem_annihilator`).
% Neither has a `\lean{...}` pin in this chapter (the lean-vs-blueprint-checker
% iter185-auslander minor flagged adding a pin for the public-facing
% `exists_isRegular_of_regularLocal` itself; iter-186+ chapter touch can add
% that hint). The Stacks 00NQ Mathlib gap remains the load-bearing substrate;
% iter-186+ work paths unchanged: (a) formalise Stacks 00NQ in project
% (~300 LOC); (b) Koszul-homology bypass via `depth_eq_smallest_ext_index`;
% (c) Mathlib upstream.
\begin{corollary}

\leanok
  [Regular local rings are Cohen--Macaulay]
  \label{cor:regular_cohen_macaulay}
  \lean{RingTheory.CohenMacaulay.of_regular}
  \uses{def:cohen_macaulay_local, thm:auslander_buchsbaum,
    lem:exists_isRegular_of_regularLocal, lem:length_le_ringKrullDim_of_isRegular}
  % SOURCE: [Stacks Project], tag 00OD (lemma-regular-ring-CM)
  % (read from references/stacks-algebra.tex, L25714-L25730).
  % SOURCE QUOTE:
  % "Let \(R\) be a regular local ring and let
  %  \(x_1, \ldots, x_d\) be a minimal set of generators
  %  for the maximal ideal \(\mathfrak m\). Then
  %  \(x_1, \ldots, x_d\) is a regular sequence, and
  %  each \(R/(x_1, \ldots, x_c)\) is a regular local ring
  %  of dimension \(d - c\). In particular \(R\) is Cohen-Macaulay."
  \textit{Source: [Stacks Project], tag 00OD (lemma-regular-ring-CM).}
  Every regular Noetherian local ring is Cohen--Macaulay: if \((R, \mathfrak{m})\) is
  regular local of Krull dimension \(d\), then \(\text{depth}(R) = d = \dim(R)\).
\end{corollary}

% SOURCE QUOTE PROOF:
% "Note that \(R/x_1R\) is a Noetherian local ring of dimension \(\geq d - 1\)
%  by Lemma REF with \(x_2, \ldots, x_d\)
%  generating the maximal ideal. Hence it is a regular local ring by definition.
%  Since \(R\) is a domain by Lemma REF
%  \(x_1\) is a nonzerodivisor."
\begin{proof}
  
  
\leanok
  Let \((R, \mathfrak{m})\) be regular of dimension \(d\) and pick a minimal generating set
  \(x_1, \dots, x_d\) of \(\mathfrak{m}\). By Stacks tag 00NQ (lemma-regular-domain), \(R\) is
  an integral domain, so \(x_1\) is a non-zero-divisor. The quotient \(R/x_1 R\) is a Noetherian
  local ring of dimension \(\geq d - 1\) (by Krull's principal ideal theorem,
  \texttt{lemma-one-equation}) with maximal ideal generated by \(x_2, \dots, x_d\), hence
  has dimension exactly \(d - 1\) and is again regular of dimension \(d - 1\). Inducting on \(d\),
  \(x_1, \dots, x_d\) is an \(R\)-regular sequence, so
  \(\text{depth}(R) \geq d\). The reverse inequality \(\text{depth}(R) \leq \dim(R) = d\)
  always holds for any nonzero finite module over a Noetherian local ring
  (Stacks tag 00LK, \texttt{lemma-bound-depth}). Thus \(\text{depth}(R) = d = \dim(R)\), i.e.\
  \(R\) is Cohen--Macaulay.

  \emph{Remark on the Auslander--Buchsbaum route.} The same conclusion can be obtained from
  \cref{thm:auslander_buchsbaum} as follows: take \(M := R\) itself. Then \(M\) is
  finitely generated and projective (free of rank \(1\)), so \(\text{pd}_R(R) = 0\). The
  formula gives \(\text{depth}(R) = 0 + \text{depth}(R)\), a tautology. The non-tautological
  use of Auslander--Buchsbaum is rather: for a regular local ring of dimension \(d\), every
  finitely generated module \(M\) has \(\text{pd}_R(M) \leq d\) (Stacks tag 00O2,
  \texttt{proposition-finite-gl-dim-regular}); applying the formula to such an \(M\) then
  relates \(\text{depth}(M)\) to \(d - \text{pd}_R(M)\), which is the form most often used in
  practice. The direct regular-sequence argument above is the cleanest proof of
  \(R\) Cohen--Macaulay itself.
\end{proof}

\section{Application: codim-1 extension on a regular surface (gating A.4.a)}
\label{sec:ab_application_to_a4a}

The downstream consumer is A.4.a: extending a rational map
\(f : S \dashrightarrow A\) from a regular projective surface \(S\) to an abelian variety \(A\)
across a closed point \(x \in S\) that is missing from the domain of definition \(U \subseteq S\).
The classical mechanism (cf.\ Hartshorne, \emph{Algebraic Geometry}, III.7) is:

\begin{enumerate}
  \item The local ring \(\mathcal{O}_{S,x}\) is a regular Noetherian local ring of Krull
    dimension \(2\) (because \(S\) is regular and \(\dim S = 2\)).
  \item By \cref{cor:regular_cohen_macaulay} it is Cohen--Macaulay, so
    \(\text{depth}(\mathcal{O}_{S,x}) = 2\). Equivalently, the maximal ideal contains an
    \(\mathcal{O}_{S,x}\)-regular sequence of length \(2\) (any two transverse generators).
  \item Sections of the structure sheaf \(\mathcal{O}_S\) on \(U = S \setminus \{x\}\) extend
    uniquely across \(x\) because \(H^0_x(\mathcal{O}_S) = 0\) and
    \(H^1_x(\mathcal{O}_S) = 0\) (these local-cohomology vanishings are equivalent to depth
    \(\geq 2\) at \(x\) by Stacks tag 0AVF; see Hartshorne, III.3.5, for the surface case).
  \item For a rational map \(f : S \dashrightarrow A\) with codimension-\(\geq 2\) indeterminacy
    locus on a regular surface, the same depth-\(\geq 2\) input produces a sheaf-theoretic
    extension of \(f\) to all of \(S\): pull back any sheaf \(\mathcal{O}_A\)-affine-locally,
    obtain sections over \(U\), and extend across \(x\) by the above.
\end{enumerate}

So \cref{cor:regular_cohen_macaulay} is the gating input: the geometric A.4.a
proof can be written conditional on the algebraic Auslander--Buchsbaum corollary, and the
two halves develop independently --- as long as \cref{thm:auslander_buchsbaum} and
\cref{cor:regular_cohen_macaulay} are formalized in
\texttt{AlgebraicJacobian/Albanese/AuslanderBuchsbaum.lean} with the signatures pinned
in this chapter, A.4.a can consume them as Mathlib-style imports.

\section{Lean encoding}
\label{sec:ab_lean_encoding}

The Lean target file is the new \texttt{AlgebraicJacobian/Albanese/AuslanderBuchsbaum.lean}.
Mathlib at the project's pinned commit (\texttt{b80f227}) already exposes substantial parts
of the depth / projective-dimension API:
\begin{itemize}
  \item \texttt{Mathlib.RingTheory.Ideal.RegularSequence} provides
    \texttt{RingTheory.Sequence.IsRegular} (definition of \(M\)-regular sequences),
    matching the regular-sequence notion of \cref{def:depth}.
  \item \texttt{Mathlib.RingTheory.Depth} (if present at the pinned commit) provides
    \texttt{Module.depth} as the supremum of regular-sequence lengths; the prover
    should re-export this as \texttt{RingTheory.Module.depth} rather than re-define.
  \item \texttt{Mathlib.CategoryTheory.Abelian.ProjectiveDimension} (or
    \texttt{Mathlib.Algebra.Homology.ProjectiveResolution}) provides
    \texttt{Module.projectiveDimension} via finite projective resolutions, matching
    \texttt{def:projective_dimension}.
  \item \texttt{Mathlib.RingTheory.CohenMacaulay} (in some recent revisions) provides
    \texttt{IsCohenMacaulayLocalRing}, matching \cref{def:cohen_macaulay_local}.
\end{itemize}
% NOTE (iter-184 review): empirical Mathlib b80f227 audit (Lane G task_result)
% confirms FOUR core ingredients of the AB-formula proof are ALL absent:
%   (i)   minimal finite free resolutions (`lemma-add-trivial-complex`);
%   (ii)  the "what is exact" criterion (Stacks 00MF);
%   (iii) snake-lemma on resolutions (tensoring minimal resolutions by R/xR);
%   (iv)  depth-drops-by-one (Stacks `lemma-depth-drops-by-one`).
% Realistic effort estimate: 4--8 iters of substrate work, NOT 30--60 LOC.
% Per the chapter's Application section, the gating consumer for A.4.a is
% Corollary~\ref{cor:regular_cohen_macaulay} (`CohenMacaulay.of_regular`), NOT
% Theorem~\ref{thm:auslander_buchsbaum} itself — AB formula is NON-BLOCKING for
% A.4.a. The plan agent may safely defer AB formula in favour of closing
% `exists_isRegular_of_regularLocal` (Stacks 00NQ; the actual `of_regular`
% blocker).
The Auslander--Buchsbaum formula itself (\cref{thm:auslander_buchsbaum}) is the
content the project must supply unless Mathlib has it at the pinned commit (the Plan Agent
should confirm via \texttt{lean\_local\_search} on \texttt{auslander\_buchsbaum} /
\texttt{Module.auslanderBuchsbaum} / similar). If absent, the file will:
\begin{enumerate}
  \item State \cref{thm:auslander_buchsbaum} as the formula
    \(\text{pd}_R(M) + \text{depth}_R(M) = \text{depth}(R)\) under the hypotheses
    ``\((R, \mathfrak{m})\) Noetherian local, \(M\) finitely generated, \(\text{pd}_R(M) < \infty\),
    \(M \neq 0\)''.
  \item Prove it by induction on \(\text{depth}(M)\) following the structure of the
    proof of \cref{thm:auslander_buchsbaum} above, importing the supporting
    \texttt{lem:depth_via_ext} from Mathlib's depth API or
    proving it in tandem.
  \item Export \cref{cor:regular_cohen_macaulay} as
    \texttt{RingTheory.CohenMacaulay.of\_regular} (or whatever name Mathlib uses if
    already present) as the consumer-facing statement.
\end{enumerate}
No new core axiom is required; the entire chapter is a homological-algebra construction.
The chapter is independent of the rest of Route A's geometric infrastructure (cube of
abelian varieties, theorem-of-the-square, Albanese universal property), making it a good
parallel-work candidate.

\section{Out of scope}
\label{sec:ab_out_of_scope}

This chapter packages \emph{only} the Auslander--Buchsbaum formula and the regular
\(\Rightarrow\) Cohen--Macaulay corollary as required by A.4.a. The following downstream
constructions are explicitly out of scope here:
\begin{itemize}
  \item \textbf{The A.4.a codim-1-to-codim-2 extension proof itself} --- consumes
    \cref{cor:regular_cohen_macaulay}, but the geometric extension argument
    lives in its sibling chapter \texttt{Albanese\_CodimExtension.tex} (planned).
  \item \textbf{Local-cohomology vanishing \(H^i_x(\mathcal{O}_S) = 0\) for \(i < \text{depth}\)}
    --- a downstream consequence of Stacks tag 0AVF; planned for its own chapter
    \texttt{Albanese\_LocalCohomology.tex} or as part of A.4.a directly.
  \item \textbf{Global Cohen--Macaulay schemes / global depth} --- this chapter only treats
    the local-ring case. The global-to-local reduction is standard but not packaged here.
  \item \textbf{General homological-dimension theory (global dimension, finitistic
    dimension)} --- only the local-ring projective-dimension invariant appears here.
  \item \textbf{Cohen--Macaulay characterisations via \(\omega_R\), dualising complex,
    Gorenstein refinement} --- the formula and its immediate Cohen--Macaulay corollary
    suffice for A.4.a; the duality side is not needed.
\end{itemize}
